\documentclass[12pt]{article}
\usepackage{amsmath, amssymb}
\usepackage{geometry}
\usepackage{enumitem}
\usepackage{array}
\geometry{margin=1in}

\title{Algebra II Final Topics Review}
\author{}
\date{}

\begin{document}
\maketitle

\section*{Instructions}
Complete in 90 minutes. Show clear work. Use exact values unless a decimal approximation is explicitly requested.

\section*{Reference Facts}
\begin{itemize}[leftmargin=1.5em]
  \item Log properties:
  \[
  \log_b(MN)=\log_bM+\log_bN,\qquad
  \log_b\!\left(\frac{M}{N}\right)=\log_bM-\log_bN,\qquad
  \log_b(M^r)=r\log_bM
  \]
  \item Hyperbola forms:
  \[
  \frac{(x-h)^2}{a^2}-\frac{(y-k)^2}{b^2}=1,\qquad
  \frac{(y-k)^2}{a^2}-\frac{(x-h)^2}{b^2}=1
  \]
  \item Arithmetic and geometric sequences:
  \[
  a_n=a_1+(n-1)d,\qquad
  a_n=a_1r^{n-1}
  \]
  \[
  S_n=\frac{n(a_1+a_n)}{2},\qquad
  S_n=\frac{a_1(1-r^n)}{1-r},\qquad
  S_\infty=\frac{a_1}{1-r}\quad(|r|<1)
  \]
  \item Counting:
  \[
  n!=n(n-1)(n-2)\cdots 1,\qquad
  {}_nP_r=\frac{n!}{(n-r)!},\qquad
  {}_nC_r=\frac{n!}{r!(n-r)!}
  \]
  \item Probability:
  \[
  P(A\text{ or }B)=P(A)+P(B)-P(A\text{ and }B),\qquad
  P(A\mid B)=\frac{P(A\text{ and }B)}{P(B)}
  \]
  \item Statistics:
  \[
  z=\frac{x-\mu}{\sigma},\qquad
  \text{range}=\max-\min,\qquad
  \mathrm{IQR}=Q_3-Q_1
  \]
\end{itemize}

\section*{Problems}
\begin{enumerate}[leftmargin=*, itemsep=1em, topsep=0.8em]
  \item Condense into a single logarithm:
  \[
  3\log_2(x)-\frac12\log_2(y+1)+2\log_2(4)-\log_2(z^3)+\log_2\!\left(\frac{w}{2}\right).
  \]

  \item Expand completely:
  \[
  \log_b\!\left(\frac{5x^3\sqrt{y}}{b^2\sqrt[3]{z^4}}\right).
  \]

  \item Solve and check for extraneous solutions:
  \[
  \log_3(x-1)+\log_3(x-7)=2+\log_3(5).
  \]

  \item Solve and state the domain restrictions used:
  \[
  \log_5(3x+2)-\log_5(x-4)=1.
  \]

  \item Given
  \[
  \log_2(a)=m,\qquad \log_2(b)=n,
  \]
  write $\log_2\!\left(\dfrac{8a^3}{\sqrt{b}}\right)$ in terms of $m$ and $n$.

  \item Determine whether each statement is always true or sometimes false. If false, give a correction or counterexample.
  \[
  \log_b(x+y)=\log_bx+\log_by,\qquad
  \log_b(x^4)=4\log_bx,\qquad
  b^{\log_b x}=x.
  \]

  \item Find the inverse function of
  \[
  f(x)=2\log_3(x-4)-5.
  \]
  State the domain of $f$ and the domain of $f^{-1}$.

  \item Graph the key features of
  \[
  y=\log_2(x+3)-4.
  \]
  State the vertical asymptote, domain, range, and $x$-intercept.

  \item A bank account earns interest compounded annually and is modeled by
  \[
  A(t)=1800(1.045)^t.
  \]
  How long will it take the account to double? Then find how long it takes to reach \$3,000. Round each to the nearest tenth of a year.

  \item Put the hyperbola in standard form and identify its center, orientation, vertices, and asymptotes:
  \[
  9x^2-16y^2-54x-64y-127=0.
  \]

  \item Write the equation of the hyperbola with center $(2,-1)$, vertices $(2,4)$ and $(2,-6)$, and asymptotes
  \[
  y+1=\pm \frac{5}{3}(x-2).
  \]
  Then find the coordinates of the foci.

  \item For the arithmetic sequence with $a_5=18$ and $S_{21}=777$, find $a_1$, $d$, and $a_{50}$.

  \item Find the exact value:
  \[
  \sum_{k=1}^{40}(7-3k).
  \]

  \item For a geometric sequence, $a_3=24$ and $a_7=\dfrac{3}{2}$. Find all possible values of $r$ and $a_1$, then find $a_{10}$ for each possible sequence.

  \item Evaluate exactly:
  \[
  \sum_{k=0}^{8}6\left(\frac{2}{3}\right)^k.
  \]

  \item Determine whether the infinite series converges. If it converges, find its exact sum:
  \[
  \sum_{n=1}^{\infty}12\left(-\frac14\right)^{n-1}.
  \]

  \item Convert the repeating decimal to a fraction:
  \[
  0.1\overline{27}.
  \]

  \item Write a recursive formula and an explicit formula for the sequence:
  \[
  7,\ 3,\ -1,\ -5,\ldots
  \]

  \item Rewrite using sigma notation:
  \[
  4+8+16+32+\cdots+512.
  \]

  \item How many different arrangements can be made from the letters in the word \textsc{LEVEL} if the two E's cannot be adjacent?

  \item A club has 9 juniors and 7 seniors. How many committees of 5 students can be formed if the committee must include at least 2 juniors and at least 2 seniors?

  \item A password must contain 2 different letters followed by 3 digits. Letters may not repeat, digits may not repeat, and the first digit cannot be 0. How many passwords are possible?

  \item A standard deck of 52 cards is shuffled and one card is drawn. Find:
  \[
  P(\text{heart or king}),\qquad
  P(\text{not a face card}).
  \]

  \item Two fair six-sided dice are rolled. Find the probability that the sum is at least 9 or the two dice show the same number.

  \item A bag contains 5 red, 4 blue, and 3 green marbles. Three marbles are drawn without replacement. Find the probability that all three are the same color.

  \item The table shows survey results from 100 students.
  \[
  \begin{array}{c|ccc}
   & \text{Likes Math} & \text{Does Not Like Math} & \text{Total}\\ \hline
  \text{Takes Art} & 18 & 22 & 40\\
  \text{Does Not Take Art} & 27 & 33 & 60\\ \hline
  \text{Total} & 45 & 55 & 100
  \end{array}
  \]
  Find $P(\text{likes math}\mid \text{takes art})$ and $P(\text{takes art}\mid \text{likes math})$. Then determine whether ``likes math'' and ``takes art'' are independent.

  \item A data set has mean $72$ and standard deviation $8$. Find the $z$-score for $x=86$. Then find the data value with $z=-1.25$.

  \item The data set is
  \[
  12,\ 15,\ 16,\ 18,\ 20,\ 24,\ 28,\ 30,\ 35.
  \]
  Find the median, $Q_1$, $Q_3$, range, and IQR.

  \item A normally distributed test has mean $80$ and standard deviation $6$. Use the empirical rule to estimate the percent of scores between $68$ and $92$, then estimate the percent above $92$.

  \item Classify each as observational study or experiment:
  \begin{enumerate}[label=(\alph*)]
    \item A researcher records how many hours students sleep and compares that to quiz scores.
    \item A researcher randomly assigns students to two study methods and compares quiz scores.
  \end{enumerate}

  \item Find the amplitude, period, midline, and phase shift of
  \[
  y=-3\sin\left(2x-\frac{\pi}{2}\right)+1.
  \]
  Then state the range.

  \item Solve on $0\le x<2\pi$:
  \[
  2\cos^2 x-\cos x-1=0.
  \]
\end{enumerate}

\end{document}
